% Complete classification of the positive equal-two-high numerator data.
% Depends on uniform_two_high_negativity_v11.tex; no geometric moduli count.

\subsection{等二高分母の整数模型の完全分類}
\label{uc11:classification-section}

前の証明の補助整数は，符号だけでなく整数模型そのものを分類する．

\begin{theorem}[等二高分母の全 numerator data]\label{uc11:classification}
正の局所四直線 chart の分母を $(q,q,1,1)$，$q>1$ とする．
二高分母の標準表示において $u=-s<0<v=z$，$q-1=tr(r+1)$ と書く．
標識の交換を除いて，可能な整数データ $(r,t,s,z)$ は次表で尽きる．
最終行では $t\geq1$ は任意の整数である．
\begin{center}
\small
\begin{tabular}{ccccccl}
$r$ & $t$ & $s$ & $z$ & $q$ & $s-z$ & $D_C+D_D$ \\
\hline
$2$ & $1$ & $2$ & $1$ & $7$ & $1$ & $2$ \\
$1$ & $6$ & $4$ & $1$ & $13$ & $3$ & $5$ \\
$2$ & $5$ & $6$ & $3$ & $31$ & $3$ & $2$ \\
$6t-3$ & $t$ & $3t$ & $3t-1$ & $t(6t-3)(6t-2)+1$ & $1$ & $3t-2$
\end{tabular}
\end{center}
ここで $C,D$ は分母一の頂点である．各行は正の weights を持ち，局所の幾何学的な chart を実現する．
従って可能な $q$ の全体は
\[
 \{7,13,31\}\ \cup\
 \{t(6t-3)(6t-2)+1:t\geq1\}
\]
である．$q=13$ では表の第二行と最終行の $t=1$ が異なる整数データを与える．
この分類は整数 numerator systems の分類であり，weights の変形や幾何学的同型類の個数を主張しない．
\end{theorem}

\begin{proof}
命題~\ref{uh11:two-integer-inequality} の証明の記号
$\delta=s-z$，$R=r(r+1)$，$n=2r+1$，$H=Z-S$，$G,J,K$ を用いる．
$\delta\geq2$ の場合，既に
\[
 \delta-1\leq J\leq3
\]
を証明したので，$\delta\leq4$ である．また同じ証明で用いた不等式を少し正確に書けば
\begin{equation}\label{uc11:sharper-J}
 J\leq\frac{\delta-1}{\delta}\left(4+\frac1R\right)
\end{equation}
となる．これは $x=z/\delta\geq r^2/n$ と
$Jx(x+1)=(\delta-1)R/\delta$ から直接従う．

\paragraph{$\delta=2,4$ の排除．}
$\delta=2$ なら \eqref{uc11:sharper-J} より $J\leq2$ である．
$J=1$ の場合は，\eqref{uh11:polynomial-elimination} の必要条件 $L>0$ に対し
$L=-8G^2+4<0$ となり矛盾する．
$J=2$ の場合は $sz=R$，$s=z+2$ だから
\[
 (z+1)^2=r^2+r+1
\]
となる．右辺は $r^2$ と $(r+1)^2$ の間に厳密にあり，平方ではない．

$\delta=4$ なら $J=3$ であり，$\delta\mid Jz$ から $z=4k$，$s=4(k+1)$ と書ける．
$Jsz=\delta(\delta-1)R$ より $R=4k(k+1)$ となるので
\[
 (2r+1)^2=4(2k+1)^2-3.
\]
従って
\[
 \{2(2k+1)-(2r+1)\}\{2(2k+1)+(2r+1)\}=3.
\]
左の二因子は正の整数なので一と三でなければならないが，その和は
$4(2k+1)\geq12$ であり矛盾する．

\paragraph{$\delta=3$ の二つの模型．}
$J=2$ または $3$ である．$J=2$ の場合には
\[
 L=-18G^2+12G+22>0
\]
から $G=1$ だけが残る．式~\eqref{uh11:polynomial-elimination} は
\[
 4R^2-23R-6=(R-6)(4R+1)=0
\]
となるので $R=6$，$r=2$ である．
$sz=18$ と $s-z=3$ から $(s,z)=(6,3)$，
$t=(\delta-1)(\delta G+J)/J$ から $t=5$ となる．

$J=3$ の場合には \eqref{uc11:sharper-J} が $R\leq2$ を与えるので $r=1$，$R=2$ である．
$sz=4$ と $s-z=3$ から $(s,z)=(4,1)$ となり，
$tK=sz+\delta-1$ と $K=1$ から $t=6$ となる．

\paragraph{$\delta=1$ の無限族と一つの例外．}
この場合は $s=z+1$，$R=HK$，$sz=tK$ であり，先の証明から $z<r$ である．
$a=r-z\geq1$ と置くと
\[
 K=R-na+1,\qquad a^2=(t-1)K+1,\qquad na=(H-1)K+1.
\]
最後の式を二乗し，$n^2=4HK+1$ と二番目の式を代入して整理すると
\begin{equation}\label{uc11:LK}
 LK=t+2H+1,\qquad L=(H+1)^2-4Ht>0
\end{equation}
となる．$L=H(H-4t+2)+1>0$ と $H\geq1$ から $H\geq4t-2$ である．
そこで $H=4t-2+k$，$k\geq0$ と書く．
\eqref{uc11:LK} は
\[
 K=\frac{9t-3+2k}{1+(4t-2)k+k^2}
\]
を与える．$k\geq3$ なら分母と分子の差は
\[
 k^2+(4t-4)k-9t+4\geq3t+1>0
\]
となり，正の整数 $K$ は存在しない．
$k=2$ なら $K=(9t+1)/(8t+1)$ は一と二の間にあり，やはり整数ではない．
$k=1$ なら $K=2+(t-1)/(4t)$ なので $t=1$ だけが可能であり，
$H=3$，$K=2$，$R=6$ から $(r,t,s,z)=(2,1,2,1)$ を得る．

残る $k=0$ の場合は
\[
 H=4t-2,\qquad K=9t-3,\qquad
 R=HK=(6t-3)(6t-2).
\]
$r>0$ なので $r=6t-3$ であり，
$a^2=(t-1)K+1=(3t-2)^2$ と $a>0$ から $a=3t-2$ となる．
従って $z=r-a=3t-1$，$s=3t$ を得る．

最後に，各行は \eqref{uh11:conic} と $s,z\mid tr(r+1)$ を満たす．
無限族では $s=3t$ の整除性は $3\mid6t-3$，$z=3t-1$ の整除性は
$6t-2=2(3t-1)$ から従い，conic は多項式恒等式である．
$u=-s<0$ の標準形なので正の weights が存在し，局所 chart の十分条件を満たす．
実際，$x=m/r$，$Z=m/z$ と書けば，全ての行で
\[
 (w_A,w_B,w_C,w_D)
 =\left(q,q,x+\frac1{2r},z+\frac1{2Z}\right)
\]
と取れる．正の numerators は $A\leftrightarrow C$ と $B\leftrightarrow D$ のみであり，
対応する四つの basin areas は $1/(2r),1/2,1/(2Z),1/2$ である．
$D_C+D_D=\delta+H-t-1$ を代入すると表の最終列を得る．
\end{proof}

無限族の高分母二頂点の係数は，標準表示の順で
\[
 D_A=2(3t-2)(t-1),\qquad
 D_B=-(3t-2)(2t-1)
\]
である．従って $t\geq2$ では一方が正である一方，その総和は常に $-(3t-2)$ である．
